\section{CLAUDE.md：一切的起点}

\begin{knowledgebox}{什么是 CLAUDE.md？}
\texttt{CLAUDE.md} 是 Claude Code 在每次对话开始时\textbf{自动读取}的指令文件，相当于给 AI 助手写的 onboarding 文档——你希望它知道什么，就写什么。不写或随便写几行的后果是：每次对话都要从头解释项目结构、编码规范和技术栈。
\end{knowledgebox}

\subsection{好的 CLAUDE.md 三原则}

\begin{importantbox}{CLAUDE.md 写作原则}
\begin{enumerate}[leftmargin=1.5em]
  \item \textbf{写代码里读不出来的东西。}项目的"为什么"比"是什么"更重要。不需要解释 React 怎么用，但要说明"选 Tailwind 是因为团队统一规范"。
  \item \textbf{控制在 200 行以内。}官方文档明确提到，CLAUDE.md 太长会导致 Claude 忽略规则。用 Markdown 标题和列表保持可扫描性。
  \item \textbf{不放频繁变化的内容。}详细的 API 文档、逐文件描述不适合放在里面，放链接就好。
\end{enumerate}
\end{importantbox}

\subsection{四级层级与规则拆分}

Claude Code 支持四级 \texttt{CLAUDE.md}，按优先级从高到低：

\begin{enumerate}[leftmargin=1.5em]
  \item \textbf{企业级}：由组织管理员统一设定
  \item \textbf{用户全局}：\texttt{\textasciitilde/.claude/CLAUDE.md}
  \item \textbf{项目级}：项目根目录下的 \texttt{CLAUDE.md}
  \item \textbf{目录级}：子目录中的 \texttt{CLAUDE.md}
\end{enumerate}

当项目变大后，可以把规则拆到 \texttt{.claude/rules/} 目录下，并用 YAML frontmatter 将规则限定到特定文件模式，这样该规则仅在 Claude 访问匹配的文件时才会加载，节省上下文空间。

\subsection{本章小结}

\texttt{CLAUDE.md} 是 Claude Code 最被低估也最值得投入的功能。写好一次，每次对话都受益。核心是：写决策理由，不写代码细节；保持精简；善用层级和规则拆分。

%% ====================================================================
